% X={1,2,3,4}上由分组生成的sigma代数：(a){1,2},{3,4}得4个集合；(b){1},{2},{3,4}得8个；(c)四个单点得16个。
\begin{tikzpicture}[x=1cm,y=1cm,
  v/.style={circle,draw=FigureStroke,fill=white,line width=.7pt,minimum size=5.5mm,inner sep=0pt,font=\fontsize{8}{9}\selectfont},
  gb/.style={draw=FigureBlue,fill=FigureBlueFill,line width=.7pt,rounded corners=5pt},
  gg/.style={draw=FigureGreen,fill=FigureGreenFill,line width=.7pt,rounded corners=5pt},
  gy/.style={draw=FigureYellow,fill=FigureYellowFill,line width=.7pt,rounded corners=5pt},
  gr/.style={draw=FigurePurple,fill=FigurePurpleFill,line width=.7pt,rounded corners=5pt},
  note/.style={font=\FigureNoteFont\fontsize{7.5}{9}\selectfont,text=FigureMuted}]
\newcommand{\pts}{\node[v] at (0,.6) {$1$};\node[v] at (1,.6) {$2$};\node[v] at (0,-.4) {$3$};\node[v] at (1,-.4) {$4$};}
\begin{scope}
  \draw[gb] (-.42,.22) rectangle (1.42,.98);
  \draw[gg] (-.42,-.78) rectangle (1.42,-.02);
  \pts
  \node[note] at (.5,-1.15) {$2^2=4$个可测集合};
  \node at (.5,-1.7) {(a) 两组};
\end{scope}
\begin{scope}[xshift=3.6cm]
  \draw[gb] (-.38,.22) rectangle (.38,.98);
  \draw[gy] (.62,.22) rectangle (1.38,.98);
  \draw[gg] (-.42,-.78) rectangle (1.42,-.02);
  \pts
  \node[note] at (.5,-1.15) {$2^3=8$个可测集合};
  \node at (.5,-1.7) {(b) 三组};
\end{scope}
\begin{scope}[xshift=7.2cm]
  \draw[gb] (-.38,.22) rectangle (.38,.98);
  \draw[gy] (.62,.22) rectangle (1.38,.98);
  \draw[gg] (-.38,-.78) rectangle (.38,-.02);
  \draw[gr] (.62,-.78) rectangle (1.38,-.02);
  \pts
  \node[note] at (.5,-1.15) {$2^4=16$个可测集合};
  \node at (.5,-1.7) {(c) 四组};
\end{scope}
\end{tikzpicture}
